\documentclass{amsart}
% For dealing with references we use the comment environment
\usepackage{verbatim}
\newenvironment{reference}{\comment}{\endcomment}
%\newenvironment{reference}{}{}
\newenvironment{slogan}{\comment}{\endcomment}
\newenvironment{history}{\comment}{\endcomment}

% For commutative diagrams we use Xy-pic
\usepackage[all]{xy}

% We use 2cell for 2-commutative diagrams.
\xyoption{2cell}
\UseAllTwocells

% We use multicol for the list of chapters between chapters
\usepackage{multicol}

% This is generally recommended for better output
\usepackage{lmodern}
\usepackage[T1]{fontenc}

% For cross-file-references

% Package for hypertext links:
\usepackage{hyperref}

% For any local file, say "hello.tex" you want to link to please

% Theorem environments.
%
\theoremstyle{plain}
\newtheorem{theorem}[subsection]{Theorem}
\newtheorem{proposition}[subsection]{Proposition}
\newtheorem{lemma}[subsection]{Lemma}

\theoremstyle{definition}
\newtheorem{definition}[subsection]{Definition}
\newtheorem{example}[subsection]{Example}
\newtheorem{exercise}[subsection]{Exercise}
\newtheorem{situation}[subsection]{Situation}

\theoremstyle{remark}
\newtheorem{remark}[subsection]{Remark}
\newtheorem{remarks}[subsection]{Remarks}

\numberwithin{equation}{subsection}

% Macros
%
\def\lim{\mathop{\mathrm{lim}}\nolimits}
\def\colim{\mathop{\mathrm{colim}}\nolimits}
\def\Spec{\mathop{\mathrm{Spec}}}
\def\Hom{\mathop{\mathrm{Hom}}\nolimits}
\def\Ext{\mathop{\mathrm{Ext}}\nolimits}
\def\SheafHom{\mathop{\mathcal{H}\!\mathit{om}}\nolimits}
\def\SheafExt{\mathop{\mathcal{E}\!\mathit{xt}}\nolimits}
\def\Sch{\mathit{Sch}}
\def\Mor{\mathop{\mathrm{Mor}}\nolimits}
\def\Ob{\mathop{\mathrm{Ob}}\nolimits}
\def\Sh{\mathop{\mathit{Sh}}\nolimits}
\def\NL{\mathop{N\!L}\nolimits}
\def\CH{\mathop{\mathrm{CH}}\nolimits}
\def\proetale{{pro\text{-}\acute{e}tale}}
\def\etale{{\acute{e}tale}}
\def\QCoh{\mathit{QCoh}}
\def\Ker{\mathop{\mathrm{Ker}}}
\def\Im{\mathop{\mathrm{Im}}}
\def\Coker{\mathop{\mathrm{Coker}}}
\def\Coim{\mathop{\mathrm{Coim}}}

% Boxtimes
%
\DeclareMathSymbol{\boxtimes}{\mathbin}{AMSa}{"02}

%
% Macros for moduli stacks/spaces
%
\def\QCohstack{\mathcal{QC}\!\mathit{oh}}
\def\Cohstack{\mathcal{C}\!\mathit{oh}}
\def\Spacesstack{\mathcal{S}\!\mathit{paces}}
\def\Quotfunctor{\mathrm{Quot}}
\def\Hilbfunctor{\mathrm{Hilb}}
\def\Curvesstack{\mathcal{C}\!\mathit{urves}}
\def\Polarizedstack{\mathcal{P}\!\mathit{olarized}}
\def\Complexesstack{\mathcal{C}\!\mathit{omplexes}}
% \Pic is the operator that assigns to X its picard group, usage \Pic(X)
% \Picardstack_{X/B} denotes the Picard stack of X over B
% \Picardfunctor_{X/B} denotes the Picard functor of X over B
\def\Pic{\mathop{\mathrm{Pic}}\nolimits}
\def\Picardstack{\mathcal{P}\!\mathit{ic}}
\def\Picardfunctor{\mathrm{Pic}}
\def\Deformationcategory{\mathcal{D}\!\mathit{ef}}

\setlength{\emergencystretch}{3em}
\begin{document}
\title[Stacks Project, Tag 03BM]{1742 --- An affine finite-locally-free equivalence relation has a finite-locally-free invariant-ring fppf quotient and is its tensor-square relation}
\maketitle
\noindent Copyright (C) 2005--2025 Johan de Jong. Stacks Project authors. Source: \href{https://stacks.math.columbia.edu/tag/03BM}{Tag 03BM}.

\noindent Editorial difficulty note (not author text). Difficulty H2: The finite-flat equivalence relation is recovered as the tensor square over its invariant ring, and the affine fppf quotient is proved finite locally free without Noetherian assumptions. The proof combines invariant norm polynomials, flat extension to infinite residue fields, the orbit description of invariant fibres, a semilocal basis chosen inside the image of one structure map, and a groupoid equalizer/basis diagram. This quotient-effectiveness and finite-projectivity construction is well beyond ordinary ring invariant calculations..

\noindent Editorial provenance note (not author text). History: excerpt from revision \texttt{a0444\allowbreak 6e57e\allowbreak c1fbc\allowbreak 252a8\allowbreak 71afc\allowbreak ec775\allowbreak 2fb28\allowbreak 07b14}, groupoids.tex, lines 4815--4910. Reference presentation was adapted; original-author context and core support blocks were retained or added. The mathematical wording is unchanged. Licensed under GNU FDL 1.2 or later; no invariant sections or cover texts. See ../licenses/COPYING. Context blocks reproduce author definitions and supporting results; they are not additional counted problems.

\noindent
Let $S$ be a scheme.
Let $(U, R, s, t, c)$ be a groupoid scheme over $S$.
Assume $U = \Spec(A)$, and $R = \Spec(B)$ are affine.
In this case we get two ring maps
$s^\sharp, t^\sharp : A \longrightarrow B$.
Let $C$ be the equalizer of $s^\sharp$ and $t^\sharp$. In a formula
\begin{equation}
\label{equation-invariants}
C = \{a \in A \mid t^\sharp(a) = s^\sharp(a) \}.
\end{equation}
We will sometimes call this the {\it ring of $R$-invariant functions on $U$}.
What properties does $M = \Spec(C)$ have? The first observation is
that the diagram
$$
\xymatrix{
R \ar[r]_s \ar[d]_t & U \ar[d] \\
U \ar[r] & M
}
$$
is commutative, i.e., the morphism $U \to M$ equalizes $s, t$.
Moreover, if $T$ is any affine scheme, and if $U \to T$ is
a morphism which equalizes $s, t$, then $U \to T$ factors through $U \to M$.
In other words, $U \to M$ is a coequalizer in the category of affine schemes.

\begin{definition}
\label{definition-quotient-sheaf}
Let $\tau$, $S$, and the pre-relation $j : R \to U \times_S U$ be as above.
In this setting the {\it quotient sheaf $U/R$} associated
to $j$ is the sheafification of the presheaf
(\ref{equation-quotient-presheaf}) in the $\tau$-topology.
If $j : R \to U \times_S U$ comes from the action of a group scheme
$G/S$ on $U$ as in Lemma \href{https://stacks.math.columbia.edu/tag/0234}{lemma groupoid from action (Tag 0234)} then we
sometimes denote the quotient sheaf $U/G$.
\end{definition}

\begin{definition}
\label{definition-groupoid}
Let $S$ be a scheme.
\begin{enumerate}
\item A {\it groupoid scheme over $S$}, or simply a
{\it groupoid over $S$} is a
quintuple $(U, R, s, t, c)$ where
$U$ and $R$ are schemes over $S$, and
$s, t : R \to U$ and $c : R \times_{s, U, t} R \to R$
are morphisms of schemes over $S$ with the
following property: For any scheme
$T$ over $S$ the quintuple
$$
(U(T), R(T), s, t, c)
$$
is a groupoid category in the sense described above.
\item A {\it morphism
$f : (U, R, s, t, c) \to (U', R', s', t', c')$
of groupoid schemes over $S$} is given by morphisms
of schemes $f : U \to U'$ and $f : R \to R'$ with the
following property:  For any scheme
$T$ over $S$ the maps $f$ define a functor from the
groupoid category $(U(T), R(T), s, t, c)$ to the
groupoid category $(U'(T), R'(T), s', t', c')$.
\end{enumerate}
\end{definition}

\begin{definition}
\label{definition-equivalence-relation}
Let $S$ be a scheme. Let $U$ be a scheme over $S$.
\begin{enumerate}
\item A {\it pre-relation} on $U$ over $S$ is any morphism
of schemes $j : R \to U \times_S U$. In this case we set
$t = \text{pr}_0 \circ j$ and $s = \text{pr}_1 \circ j$, so
that $j = (t, s)$.
\item A {\it relation} on $U$ over $S$ is a monomorphism
of schemes $j : R \to U \times_S U$.
\item A {\it pre-equivalence relation} is a pre-relation
$j : R \to U \times_S U$ such that the image of
$j : R(T) \to U(T) \times U(T)$ is an equivalence relation for
all $T/S$.
\item We say a morphism $R \to U \times_S U$ of schemes is
an {\it equivalence relation on $U$ over $S$}
if and only if for every scheme $T$ over $S$ the $T$-valued
points of $R$ define an equivalence relation
on the set of $T$-valued points of $U$.
\end{enumerate}
\end{definition}

\begin{definition}
\label{categories-definition-groupoid}
A {\it groupoid} is a category where every morphism is an isomorphism.
\end{definition}

\begin{lemma}
\label{lemma-determinant-trick}
Let $S$ be a scheme. Let $(U, R, s, t, c)$ be a groupoid scheme over $S$.
Assume $U = \Spec(A)$ and $R = \Spec(B)$ are affine and
$s, t : R \to U$ finite locally free.
Let $C$ be as in (\ref{equation-invariants}).
Let $f \in A$. Then $\text{Norm}_{s^\sharp}(t^\sharp(f)) \in C$.
\end{lemma}

\begin{proof}
Consider the commutative diagram
$$
\xymatrix{
& U & \\
R \ar[d]_s \ar[ru]^t &
R \times_{s, U, t} R
\ar[l]^-{\text{pr}_0} \ar[d]^{\text{pr}_1} \ar[r]_-c &
R \ar[d]^s \ar[lu]_t \\
U & R \ar[l]_t \ar[r]^s & U
}
$$
of Lemma \ref{lemma-diagram}.
Think of $f \in \Gamma(U, \mathcal{O}_U)$. The commutativity of the
top part of the diagram shows that
$\text{pr}_0^\sharp(t^\sharp(f)) = c^\sharp(t^\sharp(f))$ as elements of
$\Gamma(R \times_{S, U, t} R, \mathcal{O})$.
Looking at the right lower cartesian square
the compatibility of the norm construction with base change shows that
$s^\sharp(\text{Norm}_{s^\sharp}(t^\sharp(f))) =
\text{Norm}_{\text{pr}_1^\sharp}(c^\sharp(t^\sharp(f)))$.
Similarly we get
$t^\sharp(\text{Norm}_{s^\sharp}(t^\sharp(f))) =
\text{Norm}_{\text{pr}_1^\sharp}(\text{pr}_0^\sharp(t^\sharp(f)))$.
Hence by the first equality of this proof we see that
$s^\sharp(\text{Norm}_{s^\sharp}(t^\sharp(f))) =
t^\sharp(\text{Norm}_{s^\sharp}(t^\sharp(f)))$ as desired.
\end{proof}


\begin{lemma}
\label{lemma-finite-locally-free-disjoint-free}
Let $S$ be a scheme. Let $(U, R, s, t, c)$ be a groupoid scheme over $S$.
Assume $s, t : R \to U$ finite locally free.
Then
$$
U = \coprod\nolimits_{r \geq 1} U_r
$$
is a disjoint union of $R$-invariant opens such that the restriction $R_r$ of
$R$ to $U_r$ has the property that $s, t : R_r \to U_r$ are finite locally
free of rank $r$.
\end{lemma}

\begin{proof}
By
Morphisms, Lemma \href{https://stacks.math.columbia.edu/tag/04MH}{lemma finite locally free (Tag 04MH)}
there exists a decomposition
$U = \coprod\nolimits_{r \geq 0} U_r$
such that $s : s^{-1}(U_r) \to U_r$ is finite locally free of rank $r$.
As $s$ is surjective we see that $U_0 = \emptyset$.
Note that $u \in U_r \Leftrightarrow$ if and only if the scheme theoretic fibre
$s^{-1}(u)$ has degree $r$ over $\kappa(u)$. Now, if $z \in R$ with $s(z) = u$
and $t(z) = u'$ then using notation as in Lemma \ref{lemma-diagram}
$$
\text{pr}_1^{-1}(z) \to \Spec(\kappa(z))
$$
is the base change of both
$s^{-1}(u) \to \Spec(\kappa(u))$ and $s^{-1}(u') \to \Spec(\kappa(u'))$
by the lemma cited. Hence $u \in U_r \Leftrightarrow u' \in U_r$,
in other words, the open subsets $U_r$ are $R$-invariant.
In particular the restriction of $R$ to $U_r$ is just
$s^{-1}(U_r)$ and $s : R_r \to U_r$ is finite locally free of rank $r$.
As $t : R_r \to U_r$ is isomorphic to $s$ by the inverse of $R_r$
we see that it has also rank $r$.
\end{proof}


\begin{lemma}
\label{lemma-integral-over-invariants}
Let $S$ be a scheme. Let $(U, R, s, t, c)$ be a groupoid scheme over $S$.
Assume $U = \Spec(A)$ and $R = \Spec(B)$ are affine and
$s, t : R \to U$ finite locally free.
Let $C \subset A$ be as in (\ref{equation-invariants}).
Then $A$ is integral over $C$.
\end{lemma}

\begin{proof}
First, by Lemma \ref{lemma-finite-locally-free-disjoint-free}
we know that $(U, R, s, t, c)$ is a disjoint union
of groupoid schemes $(U_r, R_r, s, t, c)$ such that each $s, t : R_r \to U_r$
has constant rank $r$. As $U$ is quasi-compact, we have $U_r = \emptyset$ for
almost all $r$. It suffices to prove the lemma for each $(U_r, R_r, s, t, c)$
and hence we may assume that $s, t$ are finite locally free of rank $r$.

\medskip\noindent
Assume that $s, t$ are finite locally free of rank $r$.
Let $f \in A$. Consider the element $x - f \in A[x]$, where we think
of $x$ as the coordinate on $\mathbf{A}^1$.
Since
$$
(U \times \mathbf{A}^1, R \times \mathbf{A}^1,
s \times \text{id}_{\mathbf{A}^1},
t \times \text{id}_{\mathbf{A}^1},
c \times \text{id}_{\mathbf{A}^1})
$$
is also a groupoid scheme with finite source and target, we may apply
Lemma \ref{lemma-determinant-trick} to it and we see that
$P(x) = \text{Norm}_{s^\sharp}(t^\sharp(x - f))$
is an element of $C[x]$. Because $s^\sharp : A \to B$ is finite locally
free of rank $r$ we see that $P$ is monic of degree $r$.
Moreover $P(f) = 0$ by Cayley-Hamilton
(Algebra, Lemma \href{https://stacks.math.columbia.edu/tag/00DX}{lemma charpoly (Tag 00DX)}).
More precisely, Cayley-Hamilton says the polynomial
$s^\sharp(P) \in B[x]$ annihilates the element $t^\sharp(f)$.
Since the coefficients of $P$ are in $C$ we get that
$t^\sharp(P(f)) = 0$ in $B$. Since $t$ is faithfully flat,
we get $P(f) = 0$.
\end{proof}


\begin{lemma}
\label{lemma-invariants-base-change}
Let $S$ be a scheme. Let $(U, R, s, t, c)$ be a groupoid scheme over $S$.
Assume $U = \Spec(A)$ and $R = \Spec(B)$ are affine and
$s, t : R \to U$ finite locally free. Let $C \subset A$ be as in
(\ref{equation-invariants}). Let $C \to C'$ be a ring map, and set
$U' = \Spec(A \otimes_C C')$,
$R' = \Spec(B \otimes_C C')$.
Then
\begin{enumerate}
\item The maps $s, t, c$ induce maps $s', t', c'$ such that
$(U', R', s', t', c')$ is a groupoid scheme. Let $C^1 \subset A'$
be the $R'$-invariant functions on $U'$.
\item The canonical map $\varphi : C' \to C^1$ satisfies
\begin{enumerate}
\item for every $f \in C^1$ there exists an $n > 0$ and a
polynomial $P \in C'[x]$ whose image in $C^1[x]$ is
$(x - f)^n$, and
\item for every $f \in \Ker(\varphi)$ there exists
an $n > 0$ such that $f^n = 0$.
\end{enumerate}
\item If $C \to C'$ is flat then $\varphi$ is an isomorphism.
\end{enumerate}
\end{lemma}

\begin{proof}
The proof of part (1) is omitted. Let us denote $A' = A \otimes_C C'$ and
$B' = B \otimes_C C'$. Then we have
$$
C^1
= \{a \in A' \mid (t')^\sharp(a) = (s')^\sharp(a) \}
= \{a \in A \otimes_C C' \mid t^\sharp \otimes 1(a) = s^\sharp \otimes 1(a) \}.
$$
In other words, $C^1$ is the kernel of the difference map
$(t^\sharp - s^\sharp) \otimes 1$ which is just the base change
of the $C$-linear map $t^\sharp - s^\sharp : A \to B$ by $C \to C'$.
Hence (3) follows.

\medskip\noindent
Proof of part (2)(b). Since $C \to A$ is integral
(Lemma \ref{lemma-integral-over-invariants}) and injective we see that
$\Spec(A) \to \Spec(C)$ is surjective, see
Algebra, Lemma \href{https://stacks.math.columbia.edu/tag/00GQ}{lemma integral overring surjective (Tag 00GQ)}.
Thus also $\Spec(A') \to \Spec(C')$ is surjective
as a base change of a surjective morphism
(Morphisms, Lemma \href{https://stacks.math.columbia.edu/tag/01S1}{lemma base change surjective (Tag 01S1)}).
Hence $\Spec(C^1) \to \Spec(C')$ is surjective also.
This implies (2)(b) holds for example by
Algebra, Lemma \href{https://stacks.math.columbia.edu/tag/00FL}{lemma image dense generic points (Tag 00FL)}.

\medskip\noindent
Proof of part (2)(a). By Lemma \ref{lemma-finite-locally-free-disjoint-free}
our groupoid scheme $(U, R, s, t, c)$ decomposes as a finite disjoint union
of groupoid schemes $(U_r, R_r, s, t, c)$ such that $s, t : R_r \to U_r$
are finite locally free of rank $r$. Pulling back by $U' = \Spec(C') \to U$
we obtain a similar decomposition of $U'$ and $U^1 = \Spec(C^1)$.
We will show in the next paragraph that (2)(a) holds for the corresponding
system of rings $A_r, B_r, C_r, C'_r, C^1_r$ with $n = r$.
Then given $f \in C^1$ let $P_r \in C_r[x]$ be the polynomial
whose image in $C^1_r[x]$ is the image of $(x - f)^r$.
Choosing a sufficiently divisible integer $n$ we see that
there is a polynomial $P \in C'[x]$ whose image in $C^1[x]$ is
$(x - f)^n$; namely, we take $P$ to be the unique element of
$C'[x]$ whose image in $C'_r[x]$ is $P_r^{n/r}$.

\medskip\noindent
In this paragraph we prove (2)(a) in case the ring maps
$s^\sharp, t^\sharp : A \to B$ are finite locally free of a fixed rank $r$.
Let $f \in C^1 \subset A' = A \otimes_C C'$. Choose a flat
$C$-algebra $D$ and a surjection $D \to C'$. Choose a lift
$g \in A \otimes_C D$ of $f$.
Consider the polynomial
$$
P = \text{Norm}_{s^\sharp \otimes 1}(t^\sharp \otimes 1(x - g))
$$
in $(A \otimes_C D)[x]$. By Lemma \ref{lemma-determinant-trick}
and part (3) of the current lemma the coefficients of $P$ are in $D$
(compare with the proof of Lemma \ref{lemma-integral-over-invariants}).
On the other hand, the image of $P$ in $(A \otimes_C C')[x]$ is
$(x - f)^r$ because $t^\sharp \otimes 1(x - f) = s^\sharp(x - f)$
and $s^\sharp$ is finite locally free of rank $r$.
This proves what we want with $P$ as in the statement (2)(a)
given by the image of our $P$ under the map $D[x] \to C'[x]$.
\end{proof}


\begin{lemma}
\label{lemma-points}
Let $S$ be a scheme. Let $(U, R, s, t, c)$ be a groupoid scheme over $S$.
Assume $U = \Spec(A)$ and $R = \Spec(B)$ are affine and
$s, t : R \to U$ finite locally free. Let $C \subset A$ be as in
(\ref{equation-invariants}). Then $U \to M = \Spec(C)$ has
the following properties:
\begin{enumerate}
\item the map on points $|U| \to |M|$ is surjective and
$u_0, u_1 \in |U|$ map to the same point if and only if
there exists a $r \in |R|$ with $t(r) = u_0$ and $s(r) = u_1$, in
a formula
$$
|M| = |U|/|R|
$$
\item for any algebraically closed field $k$ we have
$$
M(k) = U(k)/R(k)
$$
\end{enumerate}
\end{lemma}

\begin{proof}
Since $C \to A$ is integral (Lemma \ref{lemma-integral-over-invariants})
and injective we see that $\Spec(A) \to \Spec(C)$ is surjective, see
Algebra, Lemma \href{https://stacks.math.columbia.edu/tag/00GQ}{lemma integral overring surjective (Tag 00GQ)}.
Thus $|U| \to |M|$ is surjective.

\medskip\noindent
Let $k$ be an algebraically closed field and let $C \to k$ be a ring map.
Since surjective morphisms are preserved under base change
(Morphisms, Lemma \href{https://stacks.math.columbia.edu/tag/01S1}{lemma base change surjective (Tag 01S1)})
we see that $A \otimes_C k$ is not zero. Now $k \subset A \otimes_C k$ is a
nonzero integral extension. Hence any residue field of $A \otimes_C k$
is an algebraic extension of $k$, hence equal to $k$. Thus we see that
$U(k) \to M(k)$ is surjective.

\medskip\noindent
Let $a_0, a_1 : A \to k$ be two ring maps. If there exists a ring map
$b : B \to k$ such that $a_0 = b \circ t^\sharp$ and $a_1 = b \circ s^\sharp$
then we see that $a_0|_C = a_1|_C$ by definition. Thus the map
$U(k) \to M(k)$ equalizes the two maps $R(k) \to U(k)$.
Conversely, suppose that $a_0|_C = a_1|_C$. Let us name this algebra
map $c : C \to k$. Consider the diagram
$$
\xymatrix{
& &
B \ar@{-->}[lld] \\
k & &
A
\ar@<0.5ex>[ll]^{a_0}
\ar@<-0.5ex>[ll]_{a_1}
\ar@<1ex>[u]
\ar@<-1ex>[u] \\
& &
C \ar[u] \ar[llu]^c
}
$$
If we can construct a dotted arrow making the diagram commute, then
the proof of part (2) of the lemma is complete. Since $s : A \to B$ is finite
there exist finitely many ring maps
$b_1, \ldots, b_n : B \to k$ such that $b_i \circ s^\sharp = a_1$.
If the dotted arrow does not exist, then we see that none of the
$a'_i = b_i \circ t^\sharp$, $i = 1, \ldots, n$ is equal to $a_0$.
Hence the maximal ideals
$$
\mathfrak m'_i = \Ker(a_i' \otimes 1 : A \otimes_C k \to k)
$$
of $A \otimes_C k$ are distinct from
$\mathfrak m = \Ker(a_0 \otimes 1 : A \otimes_C k \to k)$.
By Algebra, Lemma \href{https://stacks.math.columbia.edu/tag/00DS}{lemma silly (Tag 00DS)} we would get an element
$f \in A \otimes_C k$ with $f \in \mathfrak m$, but
$f \not \in \mathfrak m_i'$ for $i = 1, \ldots, n$.
Consider the norm
$$
g = \text{Norm}_{s^\sharp \otimes 1}(t^\sharp \otimes 1(f))
\in
A \otimes_C k
$$
By Lemma \ref{lemma-determinant-trick} this lies in the invariants
$C^1 \subset A \otimes_C k$ of the base change
groupoid (base change via the map $c : C \to k$). On the one hand,
$a_1(g) \in k^*$ since
the value of $t^\sharp(f)$ at all the points (which correspond to
$b_1, \ldots, b_n$) lying over $a_1$ is
invertible (insert future reference on property determinant here).
On the other hand, since $f \in \mathfrak m$, we see that
$f$ is not a unit, hence $t^\sharp(f)$ is not a unit
(as $t^\sharp \otimes 1$ is faithfully flat),
hence its norm is not a unit (insert future reference
on property determinant here). We conclude that $C^1$ contains
an element which is not nilpotent
and not a unit. We will now show that this leads to a contradiction.
Namely, apply Lemma \ref{lemma-invariants-base-change}
to the map $c : C \to C' = k$, then
we see that the map of $k$ into the invariants $C^1$ is injective
and moreover, that for any element $x \in C^1$ there exists an integer
$n > 0$ such that $x^n \in k$. Hence every element of $C^1$ is
either a unit or nilpotent.

\medskip\noindent
We still have to finish the proof of (1). We already know that
$|U| \to |M|$ is surjective. It is clear that $|U| \to |M|$ is
$|R|$-invariant. Finally, suppose $u_0, u_1 \in U$ maps to the same
point $m \in M$. Then the induced field extensions $\kappa(u_0)/\kappa(m)$
and $\kappa(u_1)/\kappa(m)$ are algebraic (as $A$ is integral over $C$
as used above). Hence if $k$ is an algebraic closure of $\kappa(m)$
then we can find $\kappa(m)$-embeddings $\overline{u}_0 : \kappa(u_0) \to k$
and $\overline{u}_1 : \kappa(u_1) \to k$. These determine $k$-valued
points $\overline{u}_0, \overline{u}_1 \in U(k)$ mapping to the same
point of $M(k)$. By part (2) we see that there exists a point
$\overline{r} \in R(k)$ with $s(\overline{r}) = \overline{u}_0$ and
$t(\overline{r}) = \overline{u}_1$. The image $r \in R$ of $\overline{r}$
is a point with $s(r) = u_0$ and $t(r) = u_1$ as desired.
\end{proof}


\begin{lemma}
\label{lemma-basis}
Let $S$ be a scheme.
Let $(U, R, s, t, c)$ be a groupoid scheme over $S$.
Assume
\begin{enumerate}
\item $U = \Spec(A)$, and $R = \Spec(B)$ are affine, and
\item there exist elements $x_i \in A$, $i \in I$ such that
$B = \bigoplus_{i \in I} s^\sharp(A)t^\sharp(x_i)$.
\end{enumerate}
Then $A = \bigoplus_{i\in I} Cx_i$, and $B \cong A \otimes_C A$
where $C \subset A$ is the $R$-invariant
functions on $U$ as in (\ref{equation-invariants}).
\end{lemma}

\begin{proof}
During this proof we will write $s, t : A \to B$ instead of
$s^\sharp, t^\sharp$, and similarly $c : B \to B \otimes_{s, A, t} B$.
We write $p_0 : B \to B \otimes_{s, A, t} B$, $b \mapsto b \otimes 1$ and
$p_1 : B \to B \otimes_{s, A, t} B$, $b \mapsto 1 \otimes b$. By
Lemma \ref{lemma-diagram-pull}
and the definition of $C$ we have the following
commutative diagram
$$
\xymatrix{
B \otimes_{s, A, t} B &
B \ar@<-1ex>[l]_-c \ar@<1ex>[l]^-{p_0} &
A \ar[l]^t \\
B \ar[u]^{p_1} &
A \ar@<-1ex>[l]_s \ar@<1ex>[l]^t \ar[u]_s &
C \ar[u] \ar[l]
}
$$
Moreover the tow left squares are cocartesian in the category of rings, and
the top row is isomorphic to the diagram
$$
\xymatrix{
B \otimes_{t, A, t} B &
B \ar@<-1ex>[l]_-{p_1} \ar@<1ex>[l]^-{p_0} &
A \ar[l]^t
}
$$
which is an equalizer diagram according to
Descent, Lemma \href{https://stacks.math.columbia.edu/tag/023M}{lemma ff exact (Tag 023M)} because condition (2) implies
in particular that $s$ (and hence also then isomorphic arrow $t$)
is faithfully flat.
The lower row is an equalizer diagram by definition of $C$.
We can use the $x_i$ and get a commutative diagram
$$
\xymatrix{
B \otimes_{s, A, t} B &
B \ar@<-1ex>[l]_-c \ar@<1ex>[l]^-{p_0} &
A \ar[l]^t \\
\bigoplus_{i \in I} B x_i \ar[u]^{p_1} &
\bigoplus_{i \in I} A x_i \ar@<-1ex>[l]_s \ar@<1ex>[l]^t \ar[u]_s &
\bigoplus_{i \in I} C x_i \ar[u] \ar[l]
}
$$
where in the right vertical arrow we map $x_i$ to $x_i$,
in the middle vertical arrow we map $x_i$ to $t(x_i)$ and
in the left vertical arrow we map $x_i$ to
$c(t(x_i)) = t(x_i) \otimes 1 = p_0(t(x_i))$ (equality by the commutativity
of the top part of the diagram in Lemma \ref{lemma-diagram}). Then the diagram
commutes. Moreover the middle vertical arrow is an isomorphism
by assumption. Since the left two squares are cocartesian we
conclude that also the left vertical arrow is an isomorphism.
On the other hand, the horizontal rows are exact (i.e., they are
equalizers). Hence we conclude that also the right vertical arrow
is an isomorphism.
\end{proof}


\begin{lemma}
\label{lemma-diagram-pull}
Let $S$ be a scheme.
Let $(U, R, s, t, c, e, i)$ be a groupoid over $S$.
The diagram
\begin{equation}
\label{equation-pull}
\xymatrix{
R \times_{t, U, t} R
\ar@<1ex>[r]^-{\text{pr}_1} \ar@<-1ex>[r]_-{\text{pr}_0}
\ar[d]_{(\text{pr}_0, c \circ (i, 1))} &
R \ar[r]^t \ar[d]^{\text{id}_R} &
U \ar[d]^{\text{id}_U} \\
R \times_{s, U, t} R
\ar@<1ex>[r]^-c \ar@<-1ex>[r]_-{\text{pr}_0} \ar[d]_{\text{pr}_1} &
R \ar[r]^t \ar[d]^s &
U \\
R \ar@<1ex>[r]^s \ar@<-1ex>[r]_t &
U
}
\end{equation}
is commutative. The two top rows are isomorphic via the vertical maps given.
The two lower left squares are cartesian.
\end{lemma}

\begin{proof}
The commutativity of the diagram follows from the axioms of a groupoid.
Note that, in terms of groupoids, the top left vertical arrow assigns to
a pair of morphisms $(\alpha, \beta)$ with the same target, the pair
of morphisms $(\alpha, \alpha^{-1} \circ \beta)$. In any groupoid
this defines a bijection between
$\text{Arrows} \times_{t, \text{Ob}, t} \text{Arrows}$
and
$\text{Arrows} \times_{s, \text{Ob}, t} \text{Arrows}$. Hence the second
assertion of the lemma.
The last assertion follows from Lemma \ref{lemma-diagram}.
\end{proof}


\begin{lemma}
\label{lemma-diagram}
Let $S$ be a scheme.
Let $(U, R, s, t, c)$ be a groupoid over $S$.
In the commutative diagram
$$
\xymatrix{
& U & \\
R \ar[d]_s \ar[ru]^t &
R \times_{s, U, t} R
\ar[l]^-{\text{pr}_0} \ar[d]^{\text{pr}_1} \ar[r]_-c &
R \ar[d]^s \ar[lu]_t \\
U & R \ar[l]_t \ar[r]^s & U
}
$$
the two lower squares are fibre product squares.
Moreover, the triangle on top (which is really a square)
is also cartesian.
\end{lemma}

\begin{proof}
Omitted.
Exercise in the definitions and the functorial point of
view in algebraic geometry.
\end{proof}


\begin{lemma}
\label{lemma-criterion-quotient-representable}
In the situation of Definition \ref{definition-quotient-sheaf}.
Assume there is a scheme $M$, and a morphism $U \to M$ such that
\begin{enumerate}
\item the morphism $U \to M$ equalizes $s, t$,
\item the morphism $U \to M$ induces a surjection of sheaves
$h_U \to h_M$ in the $\tau$-topology, and
\item the induced map $(t, s) : R \to U \times_M U$ induces a
surjection of sheaves $h_R \to h_{U \times_M U}$ in the $\tau$-topology.
\end{enumerate}
In this case $M$ represents the quotient sheaf $U/R$.
\end{lemma}

\begin{proof}
Condition (1) says that $h_U \to h_M$ factors through $U/R$.
Condition (2) says that $U/R \to h_M$ is surjective as a map of sheaves.
Condition (3) says that $U/R \to h_M$ is injective as a map of sheaves.
Hence the lemma follows.
\end{proof}


\begin{lemma}
\label{algebra-lemma-semi-local-module-basis-in-submodule}
Let $R$ be a local ring with maximal ideal $\mathfrak m$ and
infinite residue field.
Let $R \to S$ be a ring map.
Let $M$ be an $S$-module and let $N \subset M$ be an $R$-submodule.
Assume
\begin{enumerate}
\item $S$ is semi-local and $\mathfrak mS$ is contained in the
Jacobson radical of $S$,
\item $M$ is a finite free $S$-module, and
\item $N$ generates $M$ as an $S$-module.
\end{enumerate}
Then $N$ contains an $S$-basis of $M$.
\end{lemma}

\begin{proof}
Assume $M$ is free of rank $n$. Let $I \subset S$ be the Jacobson radical.
By Nakayama's Lemma \href{https://stacks.math.columbia.edu/tag/00DV}{lemma NAK (Tag 00DV)} a sequence of elements
$m_1, \ldots, m_n$ is a basis for $M$ if and only if
$\overline{m}_i \in M/IM$ generate $M/IM$. Hence we may replace
$M$ by $M/IM$, $N$ by $N/(N \cap IM)$, $R$ by $R/\mathfrak m$,
and $S$ by $S/IS$. In this case we see that $S$ is a finite product
of fields $S = k_1 \times \ldots \times k_r$ and
$M = k_1^{\oplus n} \times \ldots \times k_r^{\oplus n}$.
The fact that $N \subset M$ generates $M$ as an $S$-module
means that there exist $x_j \in N$ such that a linear combination
$\sum a_j x_j$ with $a_j \in S$ has a nonzero component in each
factor $k_i^{\oplus n}$.
Because $R = k$ is an infinite field, this means that also
some linear combination $y = \sum c_j x_j$ with $c_j \in k$ has a
nonzero component in each factor. Hence $y \in N$ generates a
free direct summand $Sy \subset M$. By induction on $n$ the result
holds for $M/Sy$ and the submodule $\overline{N} = N/(N \cap Sy)$.
In other words there exist $\overline{y}_2, \ldots, \overline{y}_n$
in $\overline{N}$ which (freely) generate $M/Sy$. Then
$y, y_2, \ldots, y_n$ (freely) generate $M$ and we win.
\end{proof}


\begin{lemma}
\label{algebra-lemma-flat-local-given-residue-field}
Let $(R, \mathfrak m, k)$ be a local ring. Let $K/k$ be a field
extension. There exists a local ring $(R', \mathfrak m', k')$, a flat local
ring map $R \to R'$ such that $\mathfrak m' = \mathfrak mR'$ and such that
$k'$ is isomorphic to $K$ as an extension of $k$.
\end{lemma}

\begin{proof}
Suppose that $k' = k(\alpha)$ is a monogenic extension of $k$.
Then $k'$ is the residue field of a flat local extension $R \subset R'$
as in the lemma. Namely, if $\alpha$ is transcendental over $k$, then we let
$R'$ be the localization of $R[x]$ at the prime $\mathfrak mR[x]$.
If $\alpha$ is algebraic with minimal polynomial
$T^d + \sum \overline{\lambda}_iT^{d - i}$, then we let
$R' = R[T]/(T^d + \sum \lambda_i T^{d - i})$.

\medskip\noindent
Consider the collection of triples $(k', R \to R', \phi)$, where
$k \subset k' \subset K$ is a subfield,
$R \to R'$ is a local ring map as in the lemma, and
$\phi : R' \to k'$ induces an isomorphism $R'/\mathfrak mR' \cong k'$
of $k$-extensions. These form a ``big'' category $\mathcal{C}$ with morphisms
$(k_1, R_1, \phi_1) \to (k_2, R_2, \phi_2)$
given by ring maps $\psi : R_1 \to R_2$ such that
$$
\xymatrix{
R_1 \ar[d]_\psi \ar[r]_{\phi_1} & k_1 \ar[r] & K \ar@{=}[d] \\
R_2 \ar[r]^{\phi_2} & k_2 \ar[r] & K
}
$$
commutes. This implies that $k_1 \subset k_2$.

\medskip\noindent
Suppose that $I$ is a directed set, and
$((R_i, k_i, \phi_i), \psi_{ii'})$ is a system over $I$, see
Categories, Section \href{https://stacks.math.columbia.edu/tag/002Z}{section posets limits (Tag 002Z)}.
In this case we can consider
$$
R' = \colim_{i \in I} R_i
$$
This is a local ring with maximal ideal $\mathfrak mR'$, and
residue field $k' = \bigcup_{i \in I} k_i$. Moreover, the ring
map $R \to R'$ is flat as it is a colimit of flat maps (and tensor
products commute with directed colimits).
Hence we see that $(R', k', \phi')$ is an ``upper bound'' for the system.

\medskip\noindent
An almost trivial application of Zorn's Lemma would finish the proof
if $\mathcal{C}$ was a set, but it isn't.
(Actually, you can make this work by finding a reasonable bound on the
cardinals of the local rings occurring.)
To get around this problem we choose a well ordering on $K$.
For $x \in K$ we let $K(x)$ be the subfield of $K$ generated
by all elements of $K$ which are $\leq x$.
By transfinite recursion on $x \in K$ we will produce ring maps
$R \subset R(x)$ as in the lemma with residue field extension
$K(x)/k$. Moreover, by construction we will have that
$R(x)$ will contain $R(y)$ for all $y \leq x$.
Namely, if $x$ has a predecessor $x'$, then $K(x) = K(x')[x]$
and hence we can let $R(x') \subset R(x)$ be the local ring extension
constructed in the first paragraph of the proof. If $x$ does not
have a predecessor, then we first set
$R'(x) = \colim_{x' < x} R(x')$ as in the third paragraph
of the proof. The residue field of $R'(x)$ is $K'(x) = \bigcup_{x' < x} K(x')$.
Since $K(x) = K'(x)[x]$ we see that we can use the construction of the
first paragraph of the proof to produce $R'(x) \subset R(x)$.
This finishes the proof of the lemma.
\end{proof}


Let $\tau \in \{Zariski, \etale, fppf, smooth, syntomic\}$.
Let $S$ be a scheme.
Let $j : R \to U \times_S U$ be a pre-relation over $S$.
Say $U, R, S$ are objects of a $\tau$-site $\Sch_\tau$
(see Topologies, Section \href{https://stacks.math.columbia.edu/tag/020M}{section procedure (Tag 020M)}).
Then we can consider the functors
$$
h_U, h_R :
(\Sch/S)_\tau^{opp}
\longrightarrow
\textit{Sets}.
$$
These are sheaves, see
Descent, Lemma \href{https://stacks.math.columbia.edu/tag/023Q}{lemma fpqc universal effective epimorphisms (Tag 023Q)}.
The morphism $j$ induces a map $j : h_R \to h_U \times h_U$.
For each object $T \in \Ob((\Sch/S)_\tau)$
we can take the equivalence relation $\sim_T$ generated by
$j(T) : R(T) \to U(T) \times U(T)$ and consider the quotient.
Hence we get a presheaf
\begin{equation}
\label{equation-quotient-presheaf}
(\Sch/S)_\tau^{opp}
\longrightarrow
\textit{Sets}, \quad
T \longmapsto U(T)/\sim_T
\end{equation}


\noindent
Recall that a groupoid is a category in which every morphism
is an isomorphism, see
Categories, Definition \ref{categories-definition-groupoid}.
Hence a groupoid has a set of objects $\text{Ob}$,
a set of arrows $\text{Arrows}$, a {\it source} and {\it target}
map $s, t : \text{Arrows} \to \text{Ob}$, and a {\it composition law}
$c : \text{Arrows} \times_{s, \text{Ob}, t} \text{Arrows}
\to \text{Arrows}$.
These maps satisfy exactly the following axioms
\begin{enumerate}
\item (associativity) $c \circ (1, c) = c \circ (c, 1)$ as maps
$\text{Arrows} \times_{s, \text{Ob}, t}
\text{Arrows} \times_{s, \text{Ob}, t}
\text{Arrows} \to \text{Arrows}$,
\item (identity) there exists a map $e : \text{Ob} \to \text{Arrows}$
such that
\begin{enumerate}
\item $s \circ e = t \circ e = \text{id}$ as maps $\text{Ob} \to \text{Ob}$,
\item $c \circ (1, e \circ s) = c \circ (e \circ t, 1) = 1$
as maps $\text{Arrows} \to \text{Arrows}$,
\end{enumerate}
\item (inverse) there exists a map $i : \text{Arrows} \to \text{Arrows}$
such that
\begin{enumerate}
\item $s \circ i = t$, $t \circ i = s$ as maps $\text{Arrows} \to \text{Ob}$,
and
\item $c \circ (1, i) = e \circ t$ and $c \circ (i, 1) = e \circ s$
as maps $\text{Arrows} \to \text{Arrows}$.
\end{enumerate}
\end{enumerate}
If this is the case the maps $e$ and $i$ are uniquely determined and
$i$ is a bijection. Note that if $(\text{Ob}', \text{Arrows}', s', t', c')$
is a second groupoid category, then a functor
$f : (\text{Ob}, \text{Arrows}, s, t, c) \to
(\text{Ob}', \text{Arrows}', s', t', c')$
is given by a pair of set maps $f : \text{Ob} \to \text{Ob}'$ and
$f : \text{Arrows} \to \text{Arrows}'$ such that
$s' \circ f = f \circ s$, $t' \circ f = f \circ t$, and
$c' \circ (f, f) = f \circ c$. The compatibility with identity and
inverse is automatic. We will use this below.
(Warning: The compatibility with identity
has to be imposed in the case of general categories.)


\begin{proposition}
\label{proposition-finite-flat-equivalence}
Let $S$ be a scheme.
Let $(U, R, s, t, c)$ be a groupoid scheme over $S$.
Assume
\begin{enumerate}
\item $U = \Spec(A)$, and $R = \Spec(B)$ are affine,
\item $s, t : R \to U$ finite locally free, and
\item $j = (t, s)$ is an equivalence relation.
\end{enumerate}
In this case, let $C \subset A$ be as in
(\ref{equation-invariants}). Then $U \to M = \Spec(C)$
is finite locally free and $R = U \times_M U$.
Moreover, $M$ represents the quotient sheaf $U/R$
in the fppf topology (see Definition \ref{definition-quotient-sheaf}).
\end{proposition}

\begin{proof}
During this proof we use the notation $s, t : A \to B$
instead of the notation $s^\sharp, t^\sharp$.
By Lemma \ref{lemma-criterion-quotient-representable}
it suffices to show that $C \to A$ is finite locally free
and that the map
$$
t \otimes s : A \otimes_C A \longrightarrow B
$$
is an isomorphism. First, note that $j$ is a monomorphism, and
also finite (since already $s$ and $t$ are finite). Hence we see
that $j$ is a closed immersion by
Morphisms, Lemma \href{https://stacks.math.columbia.edu/tag/03BB}{lemma finite monomorphism closed (Tag 03BB)}.
Hence $A \otimes_C A \to B$ is surjective.

\medskip\noindent
We will perform base change by flat ring maps $C \to C'$ as in
Lemma \ref{lemma-invariants-base-change}, and we will use that
formation of invariants commutes with flat base change, see
part (3) of the lemma cited.
We will show below that for every prime $\mathfrak p \subset C$, there exists
a local flat ring map $C_{\mathfrak p} \to C_{\mathfrak p}'$
such that the result holds after a base change to $C_{\mathfrak p}'$.
This implies immediately
that $A \otimes_C A \to B$ is injective (use
Algebra, Lemma \href{https://stacks.math.columbia.edu/tag/00HN}{lemma characterize zero local (Tag 00HN)}).
It also implies that $C \to A$ is flat, by combining
Algebra, Lemmas \href{https://stacks.math.columbia.edu/tag/00HR}{lemma local flat ff (Tag 00HR)},
\href{https://stacks.math.columbia.edu/tag/00HT}{lemma flat localization (Tag 00HT)}, and
\href{https://stacks.math.columbia.edu/tag/00HJ}{lemma flatness descends (Tag 00HJ)}. Then since $U \to \Spec(C)$
is surjective also (Lemma \ref{lemma-points}) we conclude that $C \to A$
is faithfully flat. Then the isomorphism $B \cong A \otimes_C A$
implies that $A$ is a finitely presented $C$-module, see
Algebra, Lemma \href{https://stacks.math.columbia.edu/tag/03C4}{lemma descend properties modules (Tag 03C4)}.
Hence $A$ is finite locally free over $C$, see
Algebra, Lemma \href{https://stacks.math.columbia.edu/tag/00NX}{lemma finite projective (Tag 00NX)}.

\medskip\noindent
By Lemma \ref{lemma-finite-locally-free-disjoint-free}
we know that $A$ is a finite
product of rings $A_r$ and $B$ is a finite product of rings $B_r$
such that the groupoid scheme decomposes accordingly (see the proof
of Lemma \ref{lemma-integral-over-invariants}).
Then also $C$ is a product of rings $C_r$ and
correspondingly $C'$ decomposes as a product. Hence we may and do
assume that the ring maps $s, t : A \to B$ are finite
locally free of a fixed rank $r$.

\medskip\noindent
The local ring maps $C_{\mathfrak p} \to C_{\mathfrak p}'$ we are going
to use are any local flat ring maps such that the residue field of
$C_{\mathfrak p}'$ is infinite.
By Algebra, Lemma \ref{algebra-lemma-flat-local-given-residue-field}
such local ring maps exist.

\medskip\noindent
Assume $C$ is a local ring with maximal ideal $\mathfrak m$ and
infinite residue field, and assume that $s, t : A \to B$ is
finite locally free of constant rank $r > 0$.
Since $C \subset A$ is integral (Lemma \ref{lemma-integral-over-invariants})
all primes lying over $\mathfrak m$ are maximal, and all maximal
ideals of $A$ lie over $\mathfrak m$. Similarly for $C \subset B$.
Pick a maximal ideal $\mathfrak m'$
of $A$ lying over $\mathfrak m$ (exists by Lemma \ref{lemma-points}).
Since $t : A \to B$ is finite locally free there exist at most finitely
many maximal ideals of $B$ lying over $\mathfrak m'$. Hence we conclude
(by Lemma \ref{lemma-points} again)
that $A$ has finitely many maximal ideals, i.e.,
$A$ is semi-local. This in turn implies that $B$ is semi-local as
well. OK, and now, because $t \otimes s : A \otimes_C A \to B$ is surjective,
we can apply
Algebra, Lemma \ref{algebra-lemma-semi-local-module-basis-in-submodule}
to the ring map $C \to A$, the $A$-module $M = B$ (seen as an $A$-module
via $t$) and the $C$-submodule $s(A) \subset B$. This lemma implies that there
exist $x_1, \ldots, x_r \in A$ such that $M$ is free over $A$
on the basis $s(x_1), \ldots, s(x_r)$. Hence we conclude that $C \to A$
is finite free and $B \cong A \otimes_C A$ by applying
Lemma \ref{lemma-basis}.
\end{proof}
\end{document}
